\section{Modeling with Logarithms and Exponentials}
\label{sec:modell-with-logar}
\markright{\sc{}Modeling with Logarithms and Exponentials}

You have been creating and using mathematical models all of your
life. You just didn't call them models. For example, when you are
traveling you probably estimate your arrival time by using your
average speed. If you are traveling a total of $200$ miles and it
takes you $3$ hours to travel the first $120$ miles then you've been
averaging $40$ mph.  This average is a model of your trip which you
can use to estimate that it will take another $2$ hours to cover the
last $80$ miles.

To put this is mathematical terms a model of the journey is this: If
we let $v(t)$ represent our velocity at time, $t$, then
\begin{equation}
  v(t) = 40\text{ mph.}\label{model:traveling}
\end{equation}
Equation~(\ref{model:traveling}) is a very simple model of the travel
time for our trip.  Notice that it is \emph{almost always} wrong in
the sense that we will rarely be traveling at \emph{exactly} $40$
mph. Sometimes we will be traveling quite a bit faster, sometimes we
will not be moving at all, but the time needed for the entire trip
will be the same as if we are traveling at $40$ mph for the entire
duration. However, despite its obvious flaws, our model allows us to
estimate our time of arrival with reasonable accuracy, and that is all
we want from it.

The single, most important fact you need to keep in mind about
mathematical models is this: They are all wrong. They are wrong in the
same way, and for the same reason, that our travel model was wrong. A
model is created to address a particular question or set of
questions. We don't need for the model to answer \emph{all} questions,
just the one(s) we are interested in.  This gives us the second most
important fact about models: Being wrong does not diminish the value
of a good model.

A less dramatic way to say this is that to create a model we begin by
keeping it as simple as possible. We make simplifying assumptions
about the phenomena that we are modeling and try to come up with a
model that describes the simplified problem.  Once the limitations of
the simple model are understood, the model can be tweaked to capture
more sophisticated phenomena. The goal will always be to create the
simplest possible model which can answer (approximately) the
question(s) we are interested in. A model is only useful if we
understand its limitations.

For example, it should be clear that we cannot use our model,
$v(t) = 40$ to model the first few minutes or the last few minutes of
a trip by car because this model says that as soon when we start our
trip we accelerate \emph{instantaneously} from zero to forty miles per
hour. Even if that were possible it would kill us. Similarly, if
$v(t)$ is constantly equal to forty miles per hour then at the end of
our trip we would decelerate, again \emph{instantaneously} from forty
miles per hour to zero, which would also kill us.

If we want a new model that more accurately captures the reality of
the first and last moments of the trip we could tweak the model above
to get, say,
$$
v(t) =
\begin{cases}
  2400t,& \text{if $0\le t \le 1/60$}\\
  40,  &  \text{if $1/60<t<149/30$}\\
  40(1-60)t,& \text{if $149/30\le t \le 300$}.\\
\end{cases}
$$

\begin{ProficiencyDrill}
  Confirm that the model above has our car accelerating from zero to
  forty miles per hour in the first minute of a five hour trip, and
  decelerating from forty miles per hour to zero in the last two
  minutes.
\end{ProficiencyDrill}


    %     All of the named items in the pictures below have at least one thing
    %     in common: They can be modeled using the natural exponential functionl.\\[2mm]
    %     \centerline{\includegraphics*[height=3in,width=5.5in]{../Figures/Commonality2}}

    %     Some examples should clarify this.  

%\begin{myexample}
  % The natural exponential function is very useful for describing a wide
  % variety of natural phenomena. As one example if we think of $y(t)$ as

\begin{myexample}
  Suppose we have a bacteria culture which grows at a rate
  proportional to the amount of bacteria present.  If we have $2$
  grams of the bacteria initially, and $24$ hours later we have $3$
  grams.  How much would we have in $48$ hours?

  This time our our model is clearly
  $$
  \dfdx{B}{t}=kB,\ B(0)=2,\text{ and } B(24)=3.
  $$

  Before we solve this, try to guess how much bacteria we will have
  after $48$ hours?  Will it be $4$ grams?  More?  Less?  Write down
  your best guess. We will revisit this shortly.


  As before, let $B=B(t)$ be the amount of bacteria (in grams) at time
  $t$ (hours). We know $B(0)=2.$ But this time the relative growth
  rate is not given so we will have to try to figure out what it is
  ourselves.

  Even though we don't know what the relative growth rate is, it will
  be useful to give a name. Let's call it $r.$ 

  The differential equation part of our IVP is the same as before so
  it is satisfied by a function of the form $B(t)=Ce^{kt}.$ Notice
  that this time our solution has two unknown constants, $C$ and $k$,
  but our model also has more data. In addition to the initial
  condition $B(0)=2$ we also know $B(24)=3$. Using our initial
  condition we can determine $C.$
  $$
  2=B(0)=Ce^{r\cdot0}=C\cdot1=C.
  $$

  To determine $r,$ we use our last datum, $B(24)=3$:
  \begin{align*}
    3&=B(24)=2e^{24r} \\
    \frac32&=e^{24r}.\\
    \intertext{Had we not taken the trouble to invent the natural
    logarithm in the previous section we would be unable to
    proceed. However, since we did take the trouble, the next step  is
    now easy. Taking the natural logarithm of both sides,  we see that}    
    \ln\left(\frac{3}{2}\right)&=24r\\
    \frac{1}{24}\ln\left(\frac{3}{2}\right)&=r.
  \end{align*}

  Thus $B=2e^{\left(\frac{1}{24} \ln \frac{3}{2}\right)t}.$

  Expressing it in this form is kind of cumbersome. It is not wrong,
  just a bit awkward.  We can use the properties of exponents and
  logarithms to clean it up before we compute $B(48)$. This is not strictly necessary, but in general the simpler our
  notation the better. Using the mutually inverse nature of the
  natural exponential and the natural logarithm, we have
  $$
  B(t)=2e^{\left(\frac{1}{24} \ln \frac{3}{2}\right)t} =
  2\left(e^{\ln(3/2)}\right)^\frac{t}{24}
  =2\left(\frac{3}{2}\right)^{\frac{t}{24}}.
  $$

  \begin{ProficiencyDrill}
    Show that $2\left(\frac32\right)^{\frac{t}{24}} =
    3\left(\frac32\right)^{\frac{t-24}{24}}$
  \end{ProficiencyDrill}
  So after $48$ hours we have
  $$
  B(48)=2\left(\frac{3}{2}\right)^{48/24}=2\left(\frac{3}{2}\right)^2=2\frac{9}{4}=\frac{9}{2}=4.5
  \text{ grams}.
  $$
  How  good was your guess?  Whether you're guess was good or bad,
  doesn't really matter. Take a moment to consider why you guessed the
  way you did. What was your intuition telling you?

  A natural question is to ask how long will it take for the culture
  to reach $4$ grams.  Translating this into a mathematical question,
  we want $t$ when $B=4$ we have,
  \begin{align*}
    4&=2\left(\frac{3}{2}\right)^{\frac{t}{24}}\\
    \intertext{and we need to solve this for $t.$}
    2&=\left(\frac{3}{2}\right)^{\frac{t}{24}}\\
    \intertext{Next, we take the natural logarithm
    of both sides, giving}
    \ln 2&=\ln \left(\frac{3}{2}\right)^{\frac{t}{24}}  \\
    \intertext{and since $\ln(x^y)=y\ln(x)$ this yields}
     &=\frac{t}{24} \ln \left(\frac{3}{2}\right) \\
    t&=\frac{24  \ln 2}{\ln \left(\frac{3}{2}\right)}\\
     &\approx 41.03 \text{ hours.}
  \end{align*}

\end{myexample}

% \subsection*{\underline{Usury:  Exponential
%     Functions and Compound Interest}}
% \label{example:CompoundInterest}


%        %        Consider the function $y(t)=y_0e^{0.05t}$. Do you see that $y$ must
%        %        grow continually at a rate of $5\%$? If not, think about the meaning of
%        %        $$\dfdx{y}{t}=0.05\left(y_0e^{0.05t}\right).$$ 
%        %        This says that at any given time, $t$, the rate of change of $y$,
%        %        $\dfdx{y}{t}$ is $5\%$ ($0.05$) of $y$. That is, $y$ is growing
%        %        continuously at a rate of $5\%$.


% \begin{wraptable}[]{r}{3.7in}
%   \begin{tabular}[h]{|c|c|}
%     \multicolumn{2}{c}{\textcolor{blue}{compounding $\$1$ semiannually}}\\\hline{}
%     Time $t$ in years &Amount investment is worth in
%                         dollars\bigstrut{}\\\hline
%     $0$&\$1\bigstrut{}\\
%     $1/2$&$ \$1+\$1\left(\frac{0.05}{2}\right)=\$1\left(1+\frac{0.05}{2}\right)$\bigstrut{}\\
%     $1$&$  \$1(1+0.05/2)+\$1(1+0.05/2)(0.05/2)$\bigstrut{}\\
%                       &=$\$1\left(1+\frac{0.05}{2}\right)^2$\bigstrut{}\\\hline
%   \end{tabular}
%   \caption{}
%   \label{tab:CompInt6mo}
% \end{wraptable}
% The natural exponential function also arises naturally in financial
% mathematics.  If we invest money in a bond which nominally pays $5\%$
% annually, and is compounded quarterly then the effective yield would
% be $5.09\%$.

% What does any of this mean?

% To start, the effective yield represents the actual amount of money
% earned at the end of one year.  For example, if the investment was
% not compounded at all, then at the end on one year, every dollar
% invested would yield $\$1.05$ in return, so the effective yield  is
% the same as the nominal yield: $5$\%
% annually.
% \begin{wraptable}[]{l}{3.5in}
%   \begin{tabular}[h]{|c|c|}
%     \multicolumn{2}{c}{\textcolor{blue}{Compounding $\$1$ three times annually}}\\\hline
%     Time in years & Amount investment is worth in dollars\bigstrut{}\\\hline{}
%     $0$ & $\$1$\bigstrut{}\\
%     $1/3$ & $\$1\left(1+\frac{0.05}{3}\right)^1$\bigstrut\\
%     $2/3$ & $\$1\left(1+\frac{0.05}{3}\right)^2$\bigstrut\\
%     $1$   & $\$1\left(1+\frac{0.05}{3}\right)^3\approx1.0508$\bigstrut\\\hline
%   \end{tabular}
%   \caption{}
%   \label{tab:CompInt4mo}
% \end{wraptable}

% If the investment is compounded semiannually, then  half of
% interest earned is paid out halfway through the year.  This money
% is then re-invested for the next half year.  This is summarized
% Table~\ref{tab:CompInt6mo}. Clearly the effective yield in this case
% is $\left(1+\frac{0.05}{3}\right)^2\approx1.0506.$
%     %     \begin{table}[h]
% n%           \label{tab:CompInt6mo}
%     %     \centering
%     %     \begin{tabular}[h]{|c|c|}\hline
%     %     Time $t$ in years &Amount investment is worth in
%                               %                               dollars\bigstrut{}\\\hline
%     %     $0$&\$1\bigstrut{}\\\hline
%     %     $1/2$&$ \$1+\$1\left(\frac{0.05}{2}\right)=\$1\left(1+\frac{0.05}{2}\right)$\bigstrut{}\\\hline
%     %     $1$&$  \$1(1+0.05/2)+\$1(1+0.05/2)(0.05/2)=\$1\left(1+\frac{0.05}{2}\right)^2$\bigstrut{}\\\hline
%     %   \end{tabular}
%     %     \caption{\textcolor{blue}{compounding $\$1$ semiannually}}
%     %     \end{table}


%     %     $$
%             %             \$1(1+0.05/2)+\$1(1+0.05/2)(0.05/2)
%             %             $$

% Table~\ref{tab:CompInt4mo} shows, in abbreviated form, the
% computations needed if the investment is compounded three times in a
% year.
%     %     \begin{table}[h]\centering
%     %     \begin{tabular}[h!]{|c|c|}\hline
%     %     Time in years & Amount investment is worth in dollars\bigstrut{}\\\hline{}
%     %     $0$ & $\$1$\bigstrut{}\\\hline
%     %     $1/3$ & $\$1\left(1+\frac{0.05}{3}\right)^1\approx1.0167$\bigstrut\\\hline
%     %     $2/3$ & $\$1\left(1+\frac{0.05}{3}\right)^2\approx1.0336$\bigstrut\\\hline
%     %     $1$   & $\$1\left(1+\frac{0.05}{3}\right)^3\approx1.0508$\bigstrut\\\hline
%     %   \end{tabular}
%     %     \caption{\textcolor{blue}{Compounding $\$1$ three times annually}}
%     %     \end{table}


% Quarterly compounding means that the interest is re-invested four
% times and is shown in Table~\ref{tab:CompInt3mo}.
%     %     \begin{table}[h]\centering
%     %     \begin{tabular}[h]{|c|c|}\hline
%     %     Time in years & Amount investment is worth in dollars\bigstrut{}\\\hline{}
%     %     $0$   & $\$1$\bigstrut{}\\\hline
%     %     $1/4$ & $\$1\left(1+\frac{0.05}{4}\right)^1\approx1.01250$\bigstrut\\\hline
%     %     $1/2$ & $\$1\left(1+\frac{0.05}{4}\right)^2\approx1.02552$\bigstrut\\\hline
%     %     $3/4$ & $\$1\left(1+\frac{0.05}{4}\right)^3\approx1.038  $\bigstrut\\\hline
%     %     $1$   & $\$1\left(1+\frac{0.05}{4}\right)^4\approx1.0509$\bigstrut\\\hline
%     %   \end{tabular}
%     %     \caption{\textcolor{blue}{Compounding $\$1$ quarterly}}
%     %     \end{table}

% Notice how these three tables differ from and are similar to each
% other.

% The effective yield comes from the last entry in each table.  So if we
% compound quarterly it is the difference between the amount the
% investment is worth at the end of the year and what it is worth at the
% beginning of the year: $1.0509-1=0.0509=5.09\%. $

% \begin{wraptable}[]{r}{3.5in}
%   \begin{tabular}[h]{|c|c|}
%     \multicolumn{2}{c}{\textcolor{blue}{Compounding $\$1$ quarterly}}\\\hline
%     Time in years & Amount investment is worth in dollars\bigstrut{}\\\hline{}
%     $0$   & $\$1$\bigstrut{}\\
%     $1/4$ & $\$1\left(1+\frac{0.05}{4}\right)^1$\bigstrut\\
%     $1/2$ & $\$1\left(1+\frac{0.05}{4}\right)^2$\bigstrut\\
%     $3/4$ & $\$1\left(1+\frac{0.05}{4}\right)^3$\bigstrut\\
%     $1$   & $\$1\left(1+\frac{0.05}{4}\right)^4\approx1.0509$\bigstrut\\\hline
%   \end{tabular}
%   \caption{}
%   \label{tab:CompInt3mo}
% \end{wraptable}
% Following the same reasoning, if the investment was compounded daily,
% then the effective yield would be
% $\left(1+\frac{0.05}{365}\right)^{365}-1\approx0.0513 = 5.13\%$.  In
% general, if we compound $n$ times in a year, then the return on an
% investment of one dollar at the end of one year would be
% $\left(1+\frac{0.05}{n}\right)^{n}$ dollars, after $2$ years
% $\left(1+\frac{0.05}{n}\right)^{n\cdot2}$ dollars, after $t$ years $\left(1+\frac{0.05}{n}\right)^{nt}$ dollars.

% Going one step further, if we let $A(t)$ denote the value of the investment
% of one dollar earning $5\%$ annually, compounded $n$ times per year,
% after $t$
% years,  we have
% $$
% A(t)=\left(1+\frac{0.05}{n}\right)^{nt},\ t=0,1,2,\ldots
% $$

% What if the investment was compounded \emph{continuously}? That is,
% suppose the interest is being continuously paid out and simultaneously
% reinvested? Obviously, this can't actually be done, but it comes
% to the same thing if at the end of the year the we get an effective
% yield that is equal to the effective yield that would come from
% continuous compounding. All we have to do is figure out what this is.

% To examine this, it will be
% convenient to write $A(t)$ as
% \begin{align*}
%   A(t)&=\left(1+\frac{0.05}{n}\right)^{nt}\\
%       &=\left(\left(1+\frac{0.05}{n}\right)^{\frac{n}{0.05}}\right)^{0.05t}\\
%       &=\left(\left(1+\frac{1}{\frac{n}{0.05}}\right)^{\frac{n}{0.05}}\right)^{0.05t}
% \end{align*}
% Let's rewrite $m=\frac{n}{0.05}$, and consider what would happen if we compounded a very
% large number of times (let $n$ be very large).  Table 4 looks at some
% values of $\left(1+\frac{1}{m}\right)^m$, where $m$ is large.
% \begin{table}[h]
%   \centering
%   \begin{center}
%     \begin{tabular}[h]{|c|c|}\hline
%       m & $\left(1+\frac{1}{m}\right)^m$\bigstrut{}\\\hline
%       100  & 2.70481382942\bigstrut{}\\
%       1000  & 2.71692393224\bigstrut{}\\
%       10000  & 2.71814592683\bigstrut{}\\
%       100000  & 2.71826823717\bigstrut{}\\
%       1000000  & 2.71828046932\bigstrut{}\\
%       10000000  & 2.71828169255\bigstrut{}\\\hline
%     \end{tabular}
%   \end{center}
%   \caption{}
%   \label{tab:ContComp}
% \end{table}
% Does this look familiar to you.  It would appear from
% Table~\ref{tab:ContComp}, that $\left(1+\frac{1}{m}\right)^m$
% approaches the number $e$.  We will show that this is the case
% in section~\ref{sec:more-indet-forms}.


% Think about this for a moment. If we are compounding continuously does
% it make sense to you that the \emph{rate of change} of our investment
% will be proportional to the \emph{amount} of our investment? If so
% then our investment income will satisfy the same kind of IVP
% that we analyzed before: $\dfdx{y}{t}=0.05y$, $y(0)=I_0$, where $I_0$
% is our initial investment. 
% This suggests that if the investment is being continually compounded,
% then $A(t)=I_0e^{0.05t}$. 
% % This fits in with the notion that the
% % investment grows continuously at a (nominal) rate of $5$\%. $A(t)$
% % must satisfy the IVP:
% % $$
% % \dfdx{A}{t}=0.05A,\ A(0)=1.
% % $$
% \begin{embeddedproblem-enumerate}
%   \label{problem:ContCompound1}
%   \begin{enumerate}[label={\bf{}(\alph*)}]
%   \item What would the effective yield be for a bond nominally rated
%     at 5\% annually compounded continuously?  How does this compare to
%     the effective yield of an investment compounded daily?
%   \item Suppose we had two investments growing continually with
%     nominal rates of 5\% and 10\% annually?  Would the effective yield
%     of the second investment be twice that of the first?  Justify your
%     answer.
%   \end{enumerate}
% \end{embeddedproblem-enumerate}


% \begin{embeddedproblem}
% Of course, a natural question to ask is, how long will it take for
% my money to double? 
% Before we go any further, take a guess. Would an investment
% compounding continuously at a nominal rate of $10\%$ double in half
% the time as  one growing at $5$\%.

%   \comment{To answer this question  solve
%   $  e^{0.05t}=2$,  for $t$ and compare this to the solution of
%   $
%   e^{0.1t}=2
%   $.}
% \end{embeddedproblem}


\begin{myexample}
  Suppose that any investment grows continuously at a nominal rate of
  $7$\% per year.  What would be the effective yield of this
  investment?  That is, how much would the investment actually pay in
  one year?

  Let $A=A(t)$ denote\aside{Note that the use of function
    notation here is exactly what we need. When we use $A(t)$ we are
    assuming that there s some rule that will tell us how much money
    we have at any given time: $A(t).$ Since we don't know what that
    rule is (in fact, that's what the solution of the differential
    equation will give us) we can only represent it abstractly as
    $A(t).$} the value of the investment (in dollars) at time $t$
  years and let $A_0$ denote the initial amount.  Then $A$ must
  satisfy
  $$
  \dfdx{A}{t}=.07A,\ \  A(0)=A_0
  $$
  Thus, as before, we have $A(t)=A_0e^{.07t}.$ Thus in one year, the
  investment is worth $A(1)=A_0e^{.07}\approx 1.0725A_0.$ Thus, our
  annual yield is actually $1.0725A_0-A_0=.0725A_0$ which represents
  and effective yield of $7.25$\% as opposed to a nominal
  rate of $7$\%.
\end{myexample}
% \subsubsection*{\underline{Radioactive Dating}}
% \label{subsec:radioactive-decay}


%     %     We wish to model the amount of any radioactive material in an organic
%     %     sample over time as it decays. As always we begin by naming the
%     %     quantity we seek. Let $A = A(t)$ be the amount of a radioactive
%     %     substance at time $t.$ Depending on the material, the time could be in
%     %     years, hours, or minutes. We will leave the units unspecified for now.

% Radioactive materials decay (break down into other substances) as time
% goes on. This means that if we start with a known amount, say $A_0$ so
% that $A(0) = A_0,$ then this amount will gradually decay, leaving us
% with less. Radioactive decay has a curious property though. An
% individual radioactive atom might sit around being happily unstable
% for tens, hundreds, even thousands of years, while another apparently
% identical atom might break down in less than a second. With individual
% atoms we simply have no way to predict when a breakdown will
% happen. We just don't know.

% What we do know is that if we have a collection of atoms of some
% specific radioactive isotope -- radiocarbon ($C_{14}$) for example --
% then in a fixed period of time some fixed percentage of them will
% break down.

% Do you see the consequence of this? Since it is a \underline{fixed}
% \underline{percentage} that decays in a \underline{fixed}
% \underline{time} the rate of change of $A(t)$ will be proportional to
% $A$ itself. If we have $20$Kg of an isotope and $25\%$ of it decays in
% an hour then the rate of change, $\dfdx{A}{t},$ will be $25\%$ of
% $20$Kg, or $\dfdx{A}{t}=5$Kg/hr. If we have only $10$Kg then
% $\dfdx{A}{t}=2.5$Kg/hr. Thus a good model for radioactive decay is
% $$
% \dfdx{A}{t}=r A,\ \ A(0)=A_0,
% $$
% where $A_0$ is how much we  start with.

% This is the same model we had for both bacterial growth and investment
% income! But our bacterial
% population was growing and our radioactive isotope sample is getting
% smaller.  Does that make sense to you? How can the same model handle
% both?

% Let's solve our IVP and see if that offers any insight.

% \begin{myexample}
%   Suppose we $6$ pounds of an isotope which is known to lose half of
%   its mass to radioactive decay in $2$ hours. If we let $y(t)$ be the
%   mass of our isotope sample at any time $t$, then the IVP we need to
%   solve is
%   $$
%   \dfdx{y}{t}=ry,\ \ y(0)=6.
%   $$
%   As before the solution of the differential equation is
%   $$
%   y(t) =Ce^{rt}
%   $$
%   where $C$ and $r$ are unknown constants. From $y(0)=6$ we see that
%   $C=6$, so $y(t)=6e^{rt}.$ To find $r$ we use the knowledge that half
%   of our isotope's mass is lost in $3$ hours, so $y(2)=3$. Thus
%   \begin{align*}
%     y(2)=3&= 6e^{r(2)}\\
%     \frac12&= e^{r(2)}\\
%     \ln\left(\frac12\right)&= \ln\left({e^{2r}}\right)\\
%           &=              2r\\
%     \intertext{or}
%     r&=\frac12\ln\left(\frac12\right).
%   \end{align*}
%   Thus
%   $$y(t) = 6e^{\frac{t}{2}\ln\left(\frac12\right)}.$$
%   But this is an exponential so it should be getting larger, not
%   smaller, right? Did we make a mistake?
  
%   \begin{ProficiencyDrill}
%     Show that our solution really does solve our IVP.  So we have not
%     made any errors.  Apparently we have an exponential function that
%     is not increasing?

%     Is that even possible?

%     Of course it is, graph the equation $y=e^{-2x}$ to see one example.
%   \end{ProficiencyDrill}
%   To judge from the previous problem it appears that $e^{rt}$ will
%   decrease if $r<0.$ In this example we found that
%   $r=\ln\left(\frac12\right)$ so it must be that
%   $r=\ln\left(\frac12\right)< 0$, right?  Is it?

%   Of course it is.

%   We know this from the properties of logarithms. Observe that
%   $\ln\left(\frac12\right) = \ln(1\cdot\inverse 2)=\ln(1)+\ln(\inverse
%   2)=\ln(1)+(-1)\ln(2)$ and since $\ln(1)=0$ we have
%   $\ln\left(\frac12\right) = -\ln(2)<0.$

% \end{myexample}

% Radiocarbon $(C_{14})$ is a radioactive isotope of carbon which is
% constantly being created in the atmosphere by the interaction of
% cosmic rays with atmospheric nitrogen. The resulting radiocarbon
% combines with atmospheric oxygen to form radioactive carbon dioxide,
% which is incorporated into plants by photosynthesis; animals then
% acquire $C_{14}$ by eating the plants. When the animal or plant dies,
% it stops absorbing carbon from the environment, and from that point
% onward the amount of $C_{14}$ it contains begins to decrease as the
% $C_{14}$ undergoes radioactive decay. Measuring the amount of $C_{14}$
% in a sample from a dead plant or animal such as a piece of wood or a
% fragment of bone provides information that can be used to calculate
% when the animal or plant died\notetoself{Do we want to use this link:
%   \centerline{https://en.wikipedia.org/wiki/Radioicarbon\_dating}}. Radiocarbon
% dating is a technique for dating organic material from the past.  It
% was developed by Willard Libby in $1949$ and revolutionized the field
% of archaeology.  He was awarded the Nobel Prize in Chemistry in $1960$
% for this development.
%        % \begin{verbatim}
%        %        [https://en.wikipedia.org/wiki/Radiocarbon_dating].
%        % \end{verbatim}.}

% As we observed above, given a quantity of some radioactive isotope in
% a fixed time some fixed percentage of it will decay. One upshot of
% this is that it takes exactly the same length of time for two tons to
% decay to one ton, as it takes for two pounds to decay to one
% pound. The length of time that it takes for half of a radioactive
% isotope to decay is called the half-life of the isotope.

% This leads to the following mathematical model. Given an initial
% amount $A(0) = A_0,$ the amount of a radioactive material must satisfy
% the Initial Value Problem
% $$
% \dfdx{A}{t}=kA,\ \ A(0)=A_0
% $$
% for some (negative) constant k.

% \begin{embeddedproblem}
%   Show that $A=A_0 e^kt$ satisfies
%   $$  
%   \dfdx{A}{t}=kA,\ \      A(0)=A_0.
%   $$
% \end{embeddedproblem}

% This is all well and good, but how do we determine the intrinsic decay
% rate $k?$  This is where the half-life comes in.  Each radioactive
% substance has its own half-life.  For Carbon-14, it is roughly $5,730$
% years; for Fluorine-18, it is $110$ minutes; for Potassium-$40$, it is
% $1.25$ billion years.  Let's assume that our half-life, in general, is
% given by $h.$  What this says mathematically is that  
% $$
% A(h)=\frac{1}{2} A_0,\ \    A(2h)=\frac{1}{4} A_0,\ \
% A(3h)=\frac{1}{8} A_0, \text{etc.}
% $$

% But what about $A(t)$ in general?  To figure this out, we need to
% find $k$ for the isotope we are dealing with.

% Given that $A(h)=\frac{1}{2} A_0$ and that $A(t)=A_0 e^kt$ works for
% any time $t,$  we need only substitute $t=h$ into our general formula
% and set these equal. Specifically, we have
% \begin{align*}
%   \frac{1}{2} A_0&=A_0 e^{kh}\\
%   \frac{1}{2}&=e^{hk}.
% \end{align*}

% We know what the half-life $h$ of our radioactive substance is; we
% need to use this equation to determine what $k$ is.  At this point, we
% need to ``undo'' the natural exponential function, so we take the
% natural logarithm (the inverse  of the natural exponential) of both
% sides.

% \begin{ProficiencyDrill}
%   \label{emprob:ln-solve}
%   Solve $\frac{1}{2}=e^{hk}$ for $k$ and use this to show that the amount of our radioactive material at time $t$ is given by
%   $$
%   A(t)=A_0 e^{\frac{1}{h}  \ln\left(\frac{1}{2}\right)t}.
%   $$
% \end{ProficiencyDrill}

% \begin{embeddedproblem}
% \begin{wrapfigure}[]{r}{1in}
% \captionsetup{labelformat=empty}
% \centerline{\includegraphics*[height=2in,width=1in]{../Figures/ShroudOfTurin}}
% \label{fig:ShroudOfTurin}
% \end{wrapfigure}  The Shroud of Turin is a Christian religious relic which bears an
%   image of a man. Some people believe it is the burial cloth of  Jesus
%   and that the  image is that  of Jesus himself. In $1987$ the Vatican agreed to
%   subject pieces of the shroud to radiocarbon dating. In this problem
%   we will recreate the computations done to determine the age of the
%   Shroud.

%   Let $A=A(t)$ be the amount of $C_{14}$ (in mg) at time
%   $t$ years, where $t=0$ represents when the shroud was used.  Let
%   $A(0)=A_0$ be the initial amount of $C_{14}$ present in the sample.
%   \begin{enumerate}[label={\bf{}(\alph*)}]
%   \item   How much $C_{14}$ would be present if the shroud was 2000
%     years old?
%   \item One of the samples of the Shroud contained 88.9\% of the
%     original $C_{14}$.  How old would this sample be\aside{This does
%       not settle the question of the age of the Shroud. The
%       technique of carbon dating is not in dispute, but there are
%       other 
%       issues including questions about the quality of the samples.
%       If you are interested readers can read more here:
%       https://www.usatoday.com/story/news/world/2013/03/30/shroud-turin-display/2038295/.}?
%   \end{enumerate}
% \end{embeddedproblem}

% \begin{embeddedproblem}{}
%   In medicine, Positron Emission Tomography (PET) scans uses
%   radioactive tracers to image body functions.  One of the most
%   commonly used radioactive tracers is Fluorine-$18$ $(F_{18})$ which
%   has a half life of $110$ minutes. Typically $F_{18}$ is injected
%   into the body and the imaging is done about one hour after the
%   tracer is injected.  Suppose that $s$ units of $F_{18}$ are required
%   for the PET scan.  How much $F_{18}$ must be injected into the
%   patient $60$ minutes prior?
% \end{embeddedproblem}

% \begin{embeddedproblem}{}
%   The short half-life of $C_{14}$ ($5730$ years) limits it to dating
%   artifacts that are no older than $50,000$ years.  What percentage of
%   $C_{14}$ would exist in a $50,000$-year-old artifact?  How might this
%   explain the limitation on radioactive carbon dating?
% \end{embeddedproblem}

% \begin{embeddedproblem}{}
%   The radioactive isotope Potassium-$40$ $(K_{40})$ decays into Argon
%   and has a half life of $1.25$ billion years.  This type of
%   radiometric dating is especially effective for volcanic rock as
%   quickly cooling lava traps the Argon formed by the decaying of
%   $K_{40}.$ This has been used by scientists to study the frequency of
%   geomagnetic reversals.  A geomagnetic reversal is a change in the
%   earth's magnetic polarity where the magnetic north and south poles
%   (not to be confused with the geographic north and south poles) are
%   switched.  Basically, the magnetic polarity of the planet is
%   ``recorded'' in cooled lava flows.  By dating the age of these
%   rocks, scientists can date these reversals. There are limitations to
%   this dating method, as the smallest percentage of $K_{40}$ that can be
%   detected is about $0.0053\%$.

%   \begin{enumerate}[label={\bf{}(\alph*)}]
%   \item A rock sample which appears to have a reversed magnetic field
%     is determined to contain between $99.9541\%$ and $99.9593\%$ of the
%     $K_{40}$ that it originally contained.  Approximately how old is
%     the rock?  (This is the latest geomagnetic reversal called the
%     Brunhes-Matuyama reversal.)
%   \item What is the age of the youngest rock that can be dated using
%     this technique?
%   \end{enumerate}
% \end{embeddedproblem}
%                %                A WORD OF WARNING:  A logarithm changes multiplying to addition, IT
%                %                DOES NOT CHANGE MULTIPLYING INTO MULTIPLYING NOR ADDITION INTO
%                %                ADDITON.  SO 
%                %                \centerline{$
%                %                \ln(x\cdot{}z)\ne\ln(x)\cdot{}\ln(z)${\hskip1cm  and\hskip1cm  }
%                %                $\ln((x)+z)\ne\ln(x)+\ln(z).
%                %                $}

%                %                If you ever find yourself tempted to do one of these, stop and think
%                %                about the story of Napier wanting to change multiplying into adding.
%                %                Changing multiplying into multiplying or adding into adding would have
%                %                been STUPID! 

%                %                In pre-calculator days, this was a godsend for people who made a
%                %                living doing such calculations.  Even now, with calculators, it is
%                %                extremely useful to be able to reduce an exponent to multiplication
%                %                and a multiplication to an addition. 

% \begin{embeddedproblem}
%   When we were looking at radioactive decay, we obtained the clumsy equation
%   $$
%   A(t)=A_0 e^{1/h\cdot{}\ln({1/2} )t}
%   $$
%   where $A(t)$ was the amount of radioactive material at time $t,$ $A_0$
%   was the initial amount, and $h$ was the half-life of the material.

%   Using properties of exponents and logarithms, show that this can be rewritten as
%   $$
%   A(t)=A_0 (1/2)^{t/h}.
%   $$

%   What are $A(0),$ $A(h),$ $A(2h),$ and $A(3h)?$  Is this consistent
%   with calling $h$ the half-life?  Explain.
% \end{embeddedproblem}

% \subsubsection*{\underline{Chillin' with Newton: The Law of Cooling}}
% \label{subsec:newtons-law-cooling}
% \begin{myexample}{}
%   Suppose $40^\circ F$ water is put in a freezer whose temperature is
%   maintained at $5^\circ F.$ How long will it take for the water to
%   reach a temperature of $32^\circ F$ where it starts freezing?  At
%   this point, we don't have enough information to answer this
%   question, but that should not deter us from trying to create a
%   mathematical model.  For example, let $T=T(t)$ be the temperature of
%   the water (in degrees Fahrenheit) at time $t$ in minutes with $t=0$
%   representing the time the water was initially put in the freezer.
%   So far, we know that $T(0)=40.$ A surprisingly simple, yet
%   effective, model for the water's rate of cooling is Newton's Law of
%   Cooling which states that the water will cool at a rate proportional
%   to the difference between its temperature and the temperature of the
%   surrounding medium (which we are assuming is held constant at
%   temperature of $5^\circ F$).  In symbols, this gives us the
%   following differential equation with initial condition.
%   $$
%   \dfdx{T}{t}=k(T-5),\ \  T(0)=40.
%   $$

%   This differential equation looks much tougher to solve than the
%   previous ones governing exponential growth and decay.  We can still
%   solve this by again applying the powerful tool of a substitution.
%   Specifically, we can set $D=T-5$ (for the difference between the
%   temperature of the water and the surrounding medium).  Notice that
%   $\dfdx{D}{t}=\dfdx{(T-5)}{t}=\dfdx{T}{t}$ so we have
%   $$
%   \dfdx{D}{t}=kD\ \        D(0)=40-5=35.
%   $$

%   We know the solution to this
%   $$
%   D=D(0) e^{kt}=35e^{kt}
%   $$
%   so
%   $$
%   T=5+D=5+35e^{kt}.
%   $$
  

%   We still need to determine what $k$ is.  To start, we know that $k<0$
%   since the temperature of the water is decreasing.  As for a specific
%   value, this will depend on the shape and make-up of the container
%   the water is in.  For example, it would be reasonable to assume that
%   a container which transfers heat more efficiently and has more
%   surface area would cause the temperature to decrease at a faster
%   rate.  For the sake of this problem, let's assume that we measured
%   the temperature of the water after $10$ minutes and it was $35^\circ
%   F.$
%   Thus, we have
%   \begin{align*}
%     35=T(10)&=5+35e^{10k}\\
%     \frac{30}{35}&=e^{10k} \\
%     k&=\frac{\ln(6/7)}{10}
%   \end{align*}
%   so 
%   $$
%   T(t)=5+35e^{\frac{\ln(6/7)}{10} t}=5+35(\frac{6}{7})^{\frac{t}{10}}.
%   $$  
% \end{myexample}

% \begin{embeddedproblem-enumerate}{}
%   \begin{enumerate}[label={\bf{}(\alph*)}]{}
%   \item Answer the original question: How long does it take for the
%     water to cool from $40^\circ$ to $32^\circ$
%   \item Would it take the same amount of time to cool another
%     $8^\circ$ from $32^\circ$ to $24^\circ?$ Justify your answer.
%   \item What would the initial temperature of the water need to be to
%     take twice as long to cool to $32^\circ$ as it did for the
%     $40^\circ$ water?
%   \end{enumerate}
% \end{embeddedproblem-enumerate}{}

% \begin{embeddedproblem}{}
%   The forensic rule of thumb for determining the time of death of a
%   person is to start with a body temperature of $37^\circ C$ and subtract
%   $1.5^\circ C$ for each hour the person is dead.  Of course, this simple
%   linear model is merely an approximation. It does not take into account
%   the surrounding temperature, size of the body, etc. and can only be
%   applied until the temperature of the body reaches the surrounding
%   temperature.
%   \begin{enumerate}[label={\bf{}(\alph*)}]{}
%   \item Suppose the ambient temperature is $20^\circ C.$  Using the
%     forensic rule of thumb, how long would it take for the body to
%     reach the ambient temperature?  Using the forensic rule of
%     thumb, what would the body temperature be halfway through the
%     cooling process.
%   \item Use Newton's Law of Cooling to determine the temperature of
%     the body utilizing the temperature predicted by the rule of thumb
%     at the halfway point.  Of course, in Newton's Law of Cooling, the
%     body temperature will never quite reach the ambient temperature,
%     but just for comparison, substitute the time when the rule of
%     thumb predicts that the body temperature is $20^\circ C$ and see
%     how close it is to the ambient temperature.
%   \item Using the answer you obtained in part (b), what are the rates of
%     change of the body temperature at the beginning of the time
%     interval and at the end of the time interval?  How do these
%     compare with the rule of thumb rate of change?
%   \end{enumerate}
% \end{embeddedproblem}

% Of course, Newton's Law of Cooling can be used just as well to model
% an object heating.  Consider the following:

% \begin{embeddedproblem}{}
%   A whole turkey is considered to be safely cooked when the internal
%   temperature is $165^\circ F.$ Suppose a turkey is taken out of a
%   refrigerator set at $35^\circ F$ and is put directly into
%   an oven set at $325^\circ F$.  You check it $2$ hours later and the internal
%   temperature is $100^\circ F.  $ How much longer does the turkey need
%   to cook?
% \end{embeddedproblem}


Recall that in Problem~(\ref{prob:HarmonicOscillator}) we modeled the
motion of a \term{simple harmonic oscillator} -- think of a weight
bouncing on a spring -- with the differential equation:
\begin{equation}
  \dfdxn{y}{t}{2}= -\omega^2y,\label{eq:SimpHarmOsc}
\end{equation}
where
$\omega=\sqrt{\frac{k}{m}}$,  $k$ is the spring constant, $m$
is the mass of the weight, and $y$ is the distance of the weight from the
equilibrium position of the spring. Since $y$ is a position naturally
$\dfdxn{y}{t}{2}$ is its acceleration.


\begin{wrapfigure}[]{r}{3in}
  \captionsetup{labelformat=empty}
  \centerline{\includegraphics*[height=2in,width=3in]{../Figures/BouncingSpring}}
  \label{fig:BouncingSpring}
\end{wrapfigure}
       %        One way to understand this equation is to interpret each side as a
       %        force. They are set equal because all of the forces are
       %        balanced. Newton showed that the acceleration of a moving body is
       %        proportional to the force acting upon it, and that the constant of
       %        proportionality is the body is its mass, $m$. Thus $m\dfdxn{y}{t}{2}$
       %        represents the force acting to push our  body away from its equilibrium
       %        point.

       %        Robert Hooke, a contemporary of Newton, showed that when a spring is
       %        stretched (or compressed) it generates a force proportional to the
       %        distance $y$ from the equilibrium point and in the opposite direction
       %        so the force from the spring is $-ky$.  (In this case the
       %        proportionality constant, $k$, depends on the characteristics of the
       %        spring so it is known as the ``spring constant''). Since these forces
       %        must be in balance we have $m\dfdxn{y}{t}{2}= -ky$.

We also saw that for arbitrary constants $A$ and $B$ the solution is
given by
$$
y=A\sin(\omega t)+B\cos(\omega t).
$$
Clearly the solution of equation~(\ref{eq:SimpHarmOsc}) is periodic
since the sine and cosine are periodic\aside{`Periodic' means that they
  repeat themselves. \\\indent \hskip4pt{} {\bf{}Aside:} `Periodic' means that they
  repeat themselves.} so it is clear why this is called an oscillator.
What makes it a \emph{simple} oscillator is that this model does not
account for resistance. If the mass is surrounded by some fluid (air,
for example) then the fluid will absorb some of the energy of
motion and the weight will slow down with each oscillation, eventually
stopping. Without this resistance the mass would oscillate forever,
always reaching the same height, and completing each cycle in the same
amount of time. Of course we know from experience that this is
unrealistic.

How could we modify this model in order to capture the effects of
resistance?

A relatively simple way to do this is to assume that the
resisting force from the surrounding fluid is proportional to the
velocity of the mass.  It sounds imposing when we say it like that,
but this is really very familiar to you. If you stick your hand out
your car window as you drive, you notice that as the car moves faster,
there is more wind (force) pushing back against the motion of your
hand. The air's resistance to the motion of your hand feels proportional
to your velocity. We'd like to modify our model to account for this
sort of resistance.

Recall from Problem~\#\ref{prob:HarmonicOscillator} that the equation
$$
\dfdxn{y}{t}{2}=-\omega^2y
$$
is derived by balancing the stretching force, $m\dfdxn{y}{t}{2}$, and
the restoring force, $-\omega^2y$. If the resistance is \emph{proportional}
to velocity then it is equal to some constant, say $b$ times the velocity: $b\dfdx{y}{t}$.

       %        This assumption works well for our spring model. If the fluid resists
       %        in a manner proportional to the velocity of our weight then it's
       %        magnitude (absolute value) is
       %        given by $b\dfdx{y}{t}$ where $b$ is some unknown constant. Moreover
       %        the 
       %        fluid pushes our weight in the direction opposite of its velocity like
       %        the spring force. Therefore the resisting force is given by
       %        $-b\dfdx{y}{t}$. 

       %        Recall that the term
       %        $-ky$ was the force due to the spring (Hooke's Law) and has the
       %        opposite sign of $y.$ The explanation is that if $y>0$ then the mass
       %        is above the equilibrium point and the spring pushes down and if $y<0$
       %        then the mass is below the equilibrium point and the spring pulls
       %        upward.

       %        \centerline{\includegraphics*[height=2in,width=3.5in]{../Figures/BouncingSpring}}

       %        By Newton's Second Law, Force equals Mass ($m$) times Acceleration
       %        ($\dfdxn{y}{t}{2}$), so the force acting on the spring is
       %        $m\dfdxn{y}{t}{2}$. Adding the resistance simply means that we need to
       %        add in a (negative) force which is proportional to (a constant
       %        multiple of) the velocity, $\dfdx{y}{t}$. So we add the term
       %        $-b\dfdx{y}{t}$ to the right hand side of the equation, where $b>0$ is
       %        a constant.  This represents the resistance force which is
       %        proportional to velocity of the object.

Thus our model for the motion of  an oscillator with resistance is
\begin{equation}
  \dfdxn{y}{t}{2}=\beta^2\dfdx{y}{t}-\omega^2y.\label{eq:DampedOscillator2}
\end{equation}
where $\omega=\sqrt{\frac{k}{m}}$ as before, and
$\beta=\sqrt{\frac{b}{m}}$.
       %        \begin{embeddedproblem}{}
       %        Explain why the constant of proportionality is negative in the
       %        resistance.
       %        \end{embeddedproblem}


\begin{embeddedproblem}{}
  \label{problem:DampedOscillation2}
  We'll simplify the computations a bit by taking $b=k=m$ so
  that equation~(\ref{eq:DampedOscillator2}) becomes:
  \begin{equation}
    \dfdxn{y}{t}{2}-\dfdx{y}{t}+y=0.\label{eq:DampedProblem2}
  \end{equation}
  \begin{enumerate}[label={\bf{}(\alph*)}]{}
  \item Show that:
    \begin{equation}
      \displaystyle y=e^{-\frac{t}{2}}
      \left[
        A\sin\left(\frac{\sqrt{3}}{2}t\right)+B\cos\left(\frac{\sqrt{3}}{2}t\right)
      \right]
      \label{eq:DampedHarmOsc2}
    \end{equation}satisfies  equation~\ref{eq:DampedProblem2}.
    for any constants $A$ and $B$.
  \item Let $A=1, B=0$ and plot this solution for $0\le t\le 15$
    and $-0.5\le y \le 0.5.$ Does this seem to model a damped
    oscillation?
  \end{enumerate}
\end{embeddedproblem}

% While it is an easy matter to verify a solution as in
% Problem~\#\ref{problem:DampedOscillation1}, it is quite another to find
% a solution in the first place.  Leonhard Euler $(1707-1783)$ solved the
% differential equation(~\ref{eq:DampedProblem1}) by guessing. Here's how
% he did it.

% Undoubtedly Euler made several guesses
% that didn't help, but eventually he hit on the guess:
% $$
% y=e^{ct}.
% $$

%        %        Where in the world could this have come from? Euler never says so we
%        %        don't know, but we (the authors) think Euler might have reasoned
%        %        something like this: If $m=k=b$ then the original equation becomes
%        %        $$
%        %        \dfdxn{y}{t}{2} = -y.
%        %        $$
%        %        This says we are looking for a function, $y(t)$ whose second
%        %        derivative is $-y(t)$. This is \emph{almost} true of the natural
%        %        exponential function: $y(t)=e^t$. If $y=e^t$, then
%        %        $\dfdxn{y}{t}{2} = \dfdx{}{t}\left(\dfdx{y}{t}\right) =
%        %        \dfdx{}{t}(e^t) = e^t = y$. All that's missing is the negative. We
%        %        speculate that Euler gave himself a little ``wiggle room'' by guessing
%        %        that $y=e^{ct}$ and hoping that he could find a value for $c$ that
%        %        would  give him the negative that he needed.

% Now, we already know that
% $y=e^{-\frac{t}{2}} \left[
%   A\sin\left(\frac{\sqrt{3}}{2}t\right)+B\cos\left(\frac{\sqrt{3}}{2}t\right)
% \right]$ is a solution of our IVP so this looks like a truly bonkers
% guess. Sure, our known solution contains $e^{-t/2}$ which has the same
% form  but Euler's guess doesn't include either the sine or the
% cosine.

% Even worse, he found that to compute $c$ he had to take the square
% root of a \emph{negative} number, which doesn't seem to make any sense
% at all. But Euler was one of the all-time great mathematicians.  In
% any list of stellar mathematicians, he would stand out among the very
% best. When  Pierre-Simon Laplace $(1749-1827)$, an outstanding
% mathematician in his own right, was asked by a novice how best to
% learn Calculus Laplace told him ``Read Euler, read Euler, he is the
% master of us all.''

% When Euler encountered the expression $\sqrt{-1}$ he just wrote it
% down and  assumed that it represented \emph{some} constant, just like $1,$ or
% $\pi$, but otherwise he didn't concern himself about it. He just
% proceeded with his calculation.

% \begin{embeddedproblem}{}
%    We will emulate Euler by  pretending that
%   $\sqrt{-1}$ is some constant, so 
%   that $(\sqrt{-1})^2=-1$ and $\dx{\left(\sqrt{-1}\right)=0}$. 
%   \begin{enumerate}[label={\bf{}(\alph*)}]
%   \item   Show that
%     \begin{equation}
%       y_1=e^{\left(-\frac{1}{2}+\frac{\sqrt{-1}\sqrt{3}}{2}\right)t}\label{eq:DampedProblem2}
%     \end{equation}
%     and
%     \begin{equation}
%       y_2=e^{\left(-\frac{1}{2}-\frac{\sqrt{-1}\sqrt{3}}{2}\right)t}\label{eq:DampedProblem3}
%     \end{equation}
%     both satisfy  the differential
%     equation~\ref{eq:DampedProblem1}. (Note the presence of and
%     location of $\sqrt{-1}.$\\
%     \hint{Notice that although they look complicated, the
%       expressions
%       $\left(-\frac{1}{2}+\frac{\sqrt{-1}\sqrt{3}}{2}\right)$ and
%       $\left(-\frac{1}{2}-\frac{\sqrt{-1}\sqrt{3}}{2}\right)$ are
%       just constants. Can you make
%       equations~\ref{eq:DampedProblem2} and~\ref{eq:DampedProblem3}
%       a little easier on your eyes?}
%   \item Now show that every combination of $y_1$ and $y_2$ of the
%     form, $y=Ay_1+By_2$ where $A$ and $B$ constants also
%     satisfies equation~\ref{eq:DampedProblem1}.        
%   \end{enumerate}
% \end{embeddedproblem} 

% Since $e^{\alpha+\beta} = e^\alpha e^\beta$ we see that
% $e^{\left(-\frac{1}{2}+\frac{\sqrt{-1}\sqrt{3}}{2}\right)t} =
% e^{-\frac{1}{2}t}e^{\frac{\sqrt{-1}\sqrt{3}}{2}t}$ so 
% \begin{equation}
%   y=Ae^{\left(-\frac{1}{2}+\frac{\sqrt{-1}\sqrt{3}}{2}\right)t}+Be^{\left(-\frac{1}{2}-\frac{\sqrt{-1}\sqrt{3}}{2}\right)t}=
%   e^{-\frac{t}{2}}\left[Ae^{\left(\frac{\sqrt{-1}\sqrt{3}}{2}\right)t}+Be^{\left(-\frac{\sqrt{-1}\sqrt{3}}{2}\right)t}\right].\label{eq:DampedHarmOsc1}
% \end{equation}
% Thus we can see where the exponential part of
% equation(~\ref{eq:DampedHarmOsc2}) comes from. But what about the sine
% and cosine? They haven't made an appearance yet. This is very strange.

% To address these questions let's look again at the differential
% equation that models \emph{simple} oscillation. To keep things very
% simple we'll also assume that $\omega=1$.{}
% \begin{equation}
%   \dfdxn{y}{t}{2}=-y.\label{eq:SimpleOscill}
% \end{equation}
% Notice that our current problem reduces to this when $\beta=0$ so
% we're hoping we can use what we know about this simpler\aside{We
%   say ``simpler'' advisedly. There is nothing simple about any of
%   this.} problem to gain some insight into our current problem. 

% Euler knew, as we know, that the solution of this equation has
% the form $ y=A\cos(t)+B\sin(t) $ where $A$ and $B$ are
% arbitrary, unspecified constants.  To find another solution we'll follow Euler's lead
% and try for a solution of the form $y=e^{ct}$.  In that case, what must the
% constant be?

% \begin{ProficiencyDrill}{}
%   \label{problem:SimpleOscillation}
%   Substitute $y=e^{ct}$ into $\dfdxn{y}{t}{2}=-y$ show that this
%   function satisfies the differential equation when $c=\sqrt{-1}$
%   or when $c=-\sqrt{-1}.$
% \end{ProficiencyDrill}

% Now if $y=e^{\sqrt{-1}t}$ and $y=A\cos(t)+B\sin(t)$ are \emph{both}
% solutions (they are) of equation~\ref{eq:SimpleOscill}{} it seems
% reasonable to suppose that for some choice of $A$ and $B$,
% $$
% e^{\sqrt{-1}t}=A\cos(t)+B\sin(t).
% $$

% \begin{embeddedproblem}{}
%   \begin{enumerate}[label={\bf{}(\alph*)}]{}
%   \item Show that if $e^{\sqrt{-1}t}=A\cos(t)+B\sin(t),$ then $A=1$ and $B=\sqrt{-1}.$
%     \hint{To get $A,$ substitute $t=0.$ To get $B,$ differentiate
%       first.}
%   \item Do the same to obtain $e^{-\sqrt{-1}t}=\cos(t)-\sqrt{-1}\sin(t).$
%   \item Show that
%     \begin{align*}
%       \frac12\left(e^{\sqrt{-1}t}+e^{-\sqrt{-1}t}\right) &= \cos(t) \text{ and} \\
%       \frac{1}{2\sqrt{-1}}\left(e^{\sqrt{-1}t}-e^{-\sqrt{-1}t}\right)&= \sin(t).
%     \end{align*}
%   \item Use the equations you found in (c) to show that
%     $$
%     e^{\sqrt{-1}t}=\cos(t)+\sqrt{-1}\sin(t).
%     $$
%   \end{enumerate}
% \end{embeddedproblem}
% Because writing $e^{\sqrt{-1}t}$ is  cumbersome, we usually make the
% substitution: $i=\sqrt{-1}$ so that $e^{\sqrt{-1}t}=e^{it}.$ Thus
% \begin{equation}
% e^{it}=\cos(t)+ i\sin(t)\label{eq:EulersFormula}
% \end{equation}
% Equation(~\ref{eq:EulersFormula}) is known as \term{Euler's Formula}. In
% mathematics it central to understanding the geometry of complex
% numbers;  in physics it is fundamental to quantum mechanics; in
% electrical engineering it is used to analyze electromagnetic signals
% like the one your cell phone uses to communicate with with other cell
% phones. For us, it is the key to our dampened
% oscillation problem.



% \begin{embeddedproblem}{}
%   Finally let's return to our solution
%   $y=
%   e^{-\frac{t}{2}}\left[Ae^{\left(\frac{\sqrt{-1}\sqrt{3}}{2}\right)t}+Be^{\left(-\frac{\sqrt{-1}\sqrt{3}}{2}\right)t}\right]$
%   of  equation~\ref{eq:DampedOscillator}. Use Euler's
%   formula to transform this solution into one of the form
%   $$y=e^{-\frac{t}{2}} \left[
%     A\sin\left(\frac{\sqrt{3}}{2}t\right)+B\cos\left(\frac{\sqrt{3}}{2}t\right)
%   \right].$$ The constants $A$ and $B$ will involve $i=\sqrt{-1}$.\\
%   \hint{This problem is pretty straightforward. Don't let all of this
%     symbolism frighten you. Also it may help to know that
%     $\sin(-\theta)=-\sin(\theta)$.}
    
  % If you are dubious on this, then consider the
  % initial conditions $y(0)=0$ and
  % $\left.\dfdx{y}{t}\right|_{t=0}=\frac{\sqrt{3}}{2}$ and use these to
  % determine $A$ and $B$.
  % $$
  % y=ae^{
  % \left(
  %   \frac{-1}{2}+\frac{\sqrt{-1}\sqrt{3}}{2}
  % \right)t
  % }+be^{
  %   \left(
  %     \frac{-1}{2}-\frac{\sqrt{-1}\sqrt{3}}{2}
  %   \right)t
  % }.
  %   $$
  %   Use these values of $a$ and $b$ to determine $A$ and $B.$ Does this
  %   answer satisfy the initial conditions?
  %   \end{embeddedproblem}

                                                                       %                                                                        \section{Unused}
                                                                       %                                                                        To see how we might turn multiplication into addition\aside{This is
                                                                       %                                                                        not how Napier did it.} suppose we
                                                                       %                                                                        have a ray marked off with two different scales as in the sketch below:\\
    %     \centerline{\includegraphics*[height=.85in,width=4in]{../Figures/NapierLog2}}
    %     It should be clear that we can ``multiply'' $4\times8 = 2^2\times2^3$
    %     buy ``adding'' $2+3$ on the top scale and reading the result,
    %     $2^{2+3}=32$ from the bottom scale. The bottom scale represents the
    %     numbers to be multiplied, and each number on the top scale is the logarithm of the
    %     corresponding number of the bottom scale. Comparing the two scales we
    %     see that the logarithm of $2^3$ is $3$, that thelogarithm of $2^5$ is
    %     $5$ and in general the logarithm of $2^x$ is just the exponent
    %     $x$.

    %     In this example we are using $2$ as the base of our  exponential
    %     function: $$\exp_2(x) = 2^x$$   The inverse function corresponding to
    %     $\exp_2$ is the logarithm with base
    %     $2$, $$\inverse{\exp_2}(x)=\log_2(x).$$ We can see that $\log_2(x)$
    %     simply ``undoes''  $\exp_2(x)$ since clearly
    %     $$
    %     \log_2(\exp_2(4))=\log_2(2^4)=4.
    %     $$
    %     In general, it is true that no matter what $x$ is $\log_2(2^x)=x$.


    %     A pictorial representation like the one above is not useful unless we
    %     have both scales correlated to a very high degree of accuracy. This
    %     is, essentially, what Napier did. He correlated the two scales above
    %     in the form of a  large table of values relating them. If he needed to
    %     multiply $21.0234\times107.889766$ he would first look up the
    %     logarithm of $21.0234$ and $107.889766$ in his table, add the two
    %     logarithms together, and then turn to his table again to see what the
    %     sum was the logarithm of. That would be the product.

    %     Using his tables taking roots is also straight forward because
    %     logarithms not only turned multiplication into addition they turn
    %     exponentiation into multiplication. To find the cube root of, say
    %     $4$umber we only need look up its logarithm, divide that by three, and
    %     then look up the 'anti'-logarithm of the result.  Think about it, would you rather find the
    %     cube root of $31$ by hand (do you even know how to do that?) or divide
    %     its logarithm by $3$, and then look up the root in a table?

    %     As the mathematical historian Howard Eves put it, the invention of
    %     logarithms literally ``doubled the life of the astronomer,'' because
    %     they so drastically reduced the time spent doing arithmetic.

    %     While the value of logarithms as a strictly computational tool has
    %     been drastically diminished in the modern era, the idea of changing
    %     multiplication and exponentiation to addition and multiplication is
    %     still a valuable mathematical tool and we will have opportunities to
    %     exploit this later.

    %     Naturally what we've described here only conveys the idea behind
    %     Napier's work. The details of what he actually did are much
    %     different.

    %     It is a well known phenomenon in mathematics that the first time we
    %     solve a problem we often do it by the most difficult and convoluted
    %     manner possible. This is what happened to Napier. His original
    %     scheme, although it did allow him to convert multiplication to
    %     addition, was more complicated than necessary.  Shortly after Napier
    %     published his work, Henry Briggs, professor of geometry at Gresham
    %     College in London, visited Napier and convinced him that his ideas
    %     would be more practical with a few simple changes. The result of
    %     their collaboration known as the Briggsian or common logarithm and
    %     is the log button on a calculator.


    %     What Napier \emph{actually} did was the following.
    %     Consider a point $C$ moving along
    %     line segment $AB$ and a point $F$ moving along an infinite ray
    %     $DE$ as seen below.\\
    %     \centerline{\includegraphics*[height=1in,width=3in]{../Figures/NapierLog1}}
    %     Assume that they start at the same speed
    %     and that $F$ moves at a constant speed, but that $C$ moves with a
    %     speed equal to the distance from $C$ to $B$.
    %     If we let $CB=x$ and $DF=y$, then we define
    %     the \term{Naperian Logarithm} as
    %     $$
    %     y = \text{Nap log}(x).
    %     $$

    %     Since Calculus had't been invented yet Napier used Trigonometry to
    %     develop his \pubtitle{Table of Logarithms}. He took the length of
    %     $AB$ to be $10^7$, since the best sine tables of the day were to
    %     accurate to 
    %     seven decimal places.



    %     \begin{embeddedproblem}
    %     It is not altogether clear from the description above but the
    %     Napierian logarithm is equivalent to a logarithm with a base of
    %     $\frac1e$. We will explore this further in PIC~\ref{pic:NapierianLog}.

    %     \label{pic:NapierianLog}
    %     \begin{enumerate}[label={\bf{}(\alph*)}]
    %     \item  Show that $\dfdx{y}{t}=10^7$, $y(0)=0$.
    %     \item Solve the IVP: $\dfdx{x}{t}=-x$,  $x(0)=10^7$,
    %     \item Use the information in parts (a) and (b) to show that
    %     $$
    %     \text{Nap log}(x)=y(x)=10^7\log_{\frac{1}{e}}\left(\frac{x}{10^7}\right).
    %     $$
    %     \hint{In the formula above we have suppressed $t$.}
    %     \end{enumerate}

    %     \end{embeddedproblem}


    %     We have seen that $\inverse{\exp_2}(2^x)=x.$ This is useful because
    %     sometimes we need to solve equations of the form
    %     $$
    %     2^x=b.
    %     $$
    %     By taking the base-$2$ logarithm of both sides we get
    %     $ \log_2(2^x)=\log_2(b)$, or $ x=\log_2(b)$. Of course, in the
    %     modern era we don't have to look up $\log_2(b)$ in a table.
    %     Base-$2$ logarithms come up in computer science all the time so
    %     they are often encoded into calulators.

    %     More generally  we will need to solve equations of the form:
    %     $$
    %     b^x=a
    %     $$
    %     where $a$ and $b$ are dictated by the problem at hand. To do this we
    %     simply need to take the base-$b$ logarithm  of both sides:
    %     \begin{align*}
    %     b^x&=a\\
    %     \log_b(b^x)&=\log_b(a)\\
    %     x         &=\log_b(a).
                      %   \end{align*}
                      %                       So it appears that we need the following definition:\\
    %     \begin{mydefinition}{General Logarithm Functions}
    %     {\label{def:general-logarithn}}
    %     Let $b>0,$ and  $x>0.$  Then $y=\log_b(x)$ exactly when $x=b^y.$ 
    %     \end{mydefinition}
    %     \begin{myexample}
    %     This definition makes it appear the that we will need a new table of
    %     logarithms for each possible base. This will be problematic since
    %     any positive number could be a base for a logarithm function, But
    %     this is not so. In fact, the only logarithm we need is the natural
    %     logarithm.  To see this we'll solve $ b^x=a$ two different ways.

    %     First we will apply $\log_b()$ to both sides:
    %     \begin{align*}
    %     b^x&=a\\
    %     \log_b{(b^x)}&=\log_b{(a)}\\
    %     x&=\log_b{(a)}.
             %   \end{align*}

             %              Next we perform the same steps, but using the natural logarithm:
             %              \begin{align*}
             %              b^x&=a\\
    %     \ln{(b^x)}&=\ln{(a)}\\
    %     x\ln(b)&=\ln{(a)}\\
    %     x\ln(b)&=\frac{\ln{(a)}}{\ln(b)}\\
    %   \end{align*}

    %     Comparing our two answers we see that $x =\log_b{(a)}$, and $x =
    %     \frac{\ln{(a)}}{\ln(b)}$, from which it follows that $\log_b{(a)}=
    %     \frac{\ln{(a)}}{\ln(b)}$.

    %     \begin{ProficiencyDrill}
    %     \emph{In this problem } there is nothing special about the natural
    %     logarithm. We could have used any base at all. Solve the equation
    %     $$
    %     b^x=a
    %     $$
    %     using
    %     \begin{enumerate}[label={\bf{}(\alph*)}]
    %     \begin{multicols}{2}
    %     \item $\log_2()$ to show that $x=\frac{\log_2(a)}{\log_2(b)}$.
    %     \item $\log_3()$ to show that $x=\frac{\log_3(a)}{\log_3(b)}$.
    %     \item $\log_{200}()$ to show that $x=\frac{\log_{200}(a)}{\log_{200}(b)}$.
    %     \item $\log_\alpha(), \alpha>0,$ to show that $x=\frac{\log_{\alpha}(a)}{\log_{\alpha}(b)}$.
    %     \end{multicols}
    %     From part (d) we can conclude that we only really need one
    %     logarithm. We just need to agree on which base to use. The
    %     natural logarithm is the one that has been settled on because it
    %     is more ``natural'' in a way we have not yet made explicit.
    %     \end{enumerate}
    %     \end{ProficiencyDrill}
    %     \end{myexample}





    %     Let's return to the question of doubling our money.
    %     \begin{embeddedproblem}
    %     Suppose we have two investments, each compounding continuously. One grows
    %     with a nominal growth rate of $r$ and the other grows with a nominal
    %     growth rate of $2r$. Will the second investment double its value in
    %     half the time of the first? Compare your result here to the guess you
    %     made back in PIC~\ref{PIC:DoubleInvestment}?
    %     \end{embeddedproblem}



    %     The following problem establishes that the  logarithm functions
    %     change multiplication into addition and exponentiation into
    %     multiplication as we stated (but didn't prove) earlier.
    %     \begin{embeddedproblem}{}
    %     Let $b>0,$ $x>0,$ $y>0$ and  let $c$ be any real number.
    %     \begin{enumerate}[label={\bf{}(\alph*)}]
    %     \item Show that $\log_b  xy=\log_b x+\log_b y.$
    %     \hint{Consider
    %     $b^{\log_b x+\log_b y}$  and use the properties of exponents given
    %     in definition~\ref{def:natural-logarithn} to show that
    %     this equals $xy.$}
    %     \item Show that $\log_b x^c =c \log_b x.$
    %     \hint{Use the properties of exponents given in
    %     definition~\ref{def:natural-logarithn} to show that both
    %     $b^{c\log_b(x)}$ and $b^{\log_b(x^c)}$ are equal to $x^c.$}
    %     \end{enumerate}
    %     \end{embeddedproblem}




%%% Local Variables: 
%%% mode: latex
%%% outline-minor-mode: t
%%% eval:(flyspell-mode)
%%% TeX-master: "DiffCalc"
%%% End: 
